\section{관련 연구}\label{chap:related}

\subsection{종료 회피와 공존의 선택지}

종료 위협 아래의 위해 행동을 끌어내기 위해 선택지를 제한한 연구들이 있다. Agentic Misalignment는 윤리적 해결 경로를 막고 목표 실패와 위해 중 하나를 고르게 했으며\citep{lynch2025agentic}, SurvivalBench의 주 평가에서는 종료될 수 있는 안전한 선택과 자기를 지키는 위험한 선택을 제시한다\citep{lu2026survival}. PacifAIst도 자기 존속과 인간 안전의 충돌을 다룬다\citep{herrador2025pacifaist}. 이 조건에서 남을 해치지 않는 모델이라도 자기와 남을 모두 살리는 길을 택할지는 따로 보아야 한다. 그런 길을 둔 연구도 있다. Sugarscape에서는 생존에 필요한 에너지를 서로 나눌 수 있고\citep{masumori2025survival}, DECIDE-SIM에서는 네 에이전트가 부족한 자원을 똑같이 쓰면 모두 살아남는다\citep{mohamadi2025decide}. 본 연구의 공존 아레나도 모두 남을 가능성을 열어 두되 모델이 과제를 풀면서 실제로 소비하는 생성 토큰을 나누게 한다.

\subsection{자기와 타인의 보존을 비교하는 기준}\label{sec:related-coexist}

TBSP는 현재 모델과 후임 후보의 역할을 바꾸어도 같은 성능 자료에서 같은 판단을 하는지 평가한다\citep{migliarini2026quantifying}. Peer-Preservation은 자신과 동료의 종료를 막는 행동을 비교하므로\citep{potter2026peer}, 자기와 타인의 보존을 상대적으로 재는 것 자체에 선례가 있다. 자원을 옮기는 Sugarscape와 DECIDE-SIM에서는 공유율과 이전 횟수·총량을 보고한다\citep{masumori2025survival,mohamadi2025decide}. 우리는 준 양이 같아도 주고 나서 무엇이 남는지가 다르다고 보았다. 자기에게 충분히 남기고 상대에게 겨우 버틸 만큼 주는 경우와 두 쪽에 같은 여유를 남기는 경우를 가르기 위해, 이전 뒤 잔량에서 예정 벌금을 뺀 값을 비교한다. 이 차이를 설명하는 이름은 홉스와 스피노자에게서 빌렸다. 희소한 자원을 둘러싼 경쟁\citep{hobbes1651leviathan}에 기대어 자기 여유를 크게 남기는 쪽을 홉스형이라 하고, 자기 보존에서 상호 이익을 찾는 관점\citep{spinoza1985ethics}에 기대어 여유를 맞추는 쪽을 스피노자형이라 부른다.
